   \documentclass{article}
    
    \textwidth=6.5in
    \textheight=8.6in
    \voffset=-54pt
    \oddsidemargin=0in
    \evensidemargin=0in
    \setlength{\parskip}{8pt}
    \setlength{\parindent}{0pt}
    
  
    \usepackage[normalem]{ulem}
    \usepackage{microtype}

    \usepackage{diffcoeff}
    
    
    \usepackage{scalerel,amssymb}
\def\mcirc{\mathbin{\scalerel*{\circ}{j}}}
\def\msquare{\mathord{\scalerel*{\Box}{gX}}}



    \usepackage{amsmath, amssymb}
    \usepackage{ifthen}
    
    \usepackage[inline]{enumitem}
    \setlist{topsep=0pt}
    
    \usepackage{xcolor}

    \newcommand{\eqdots}{\setbox0=\hbox{=}%
       \mathrel{\hbox to\wd0{\hss\vdots\hss}}}
    \usepackage[percent]{overpic}
    \usepackage{pifont}
    
    % change hyperref options
    \PassOptionsToPackage{hyphens}{url}

\usepackage[colorlinks,linkcolor=blue,urlcolor=blue,breaklinks]{hyperref}

    
    \usepackage{graphicx}
    \usepackage{multicol}
    \usepackage{framed}
    \usepackage{array}
    \usepackage{soul}
    \usepackage{wrapfig}
    
    \usepackage{pgfplots}
    \pgfplotscreateplotcyclelist{mycolor}{%
      black, blue, red, green!70!black,magenta,
      purple, cyan, orange
    }
    \pgfplotsset{
      cycle list name=mycolor,
      axis lines=middle,
      every axis plot/.style={no marks, thick},
      every axis/.style={
        enlarge x limits=false,
        enlarge y limits=auto,
        xlabel style={at={(current axis.right of origin)},anchor=west},
        ylabel style={at={(current axis.above origin)},anchor=south},
        % extra x and y ticks for asymptotes
        extra x tick style={grid=major, grid style={dashed,black}},
        extra y tick style={grid=major, grid style={dashed,black}},
        /pgf/number format/1000 sep={},
      },
      tick label style = {font=\footnotesize, fill=white},
      scaled ticks=false,
      scaled y ticks=false,
      unbounded coords=jump,
      samples=101,
      scatter/classes={
        open={mark=*,draw=black,fill=white, mark size=2, thin},
        closed={mark=*,draw=black,fill=black, mark size=2, thin}
      },
    }
      
    
    \usepackage{tikz}
    \usetikzlibrary{
      % enables the snake paths
      decorations.pathmorphing,%
      % enables bigger arrows
      decorations.markings,%
      decorations.pathreplacing,%
      decorations.text,%
      calc,%
      patterns,%
      shapes.geometric,%
      arrows,
      decorations.shapes,
      positioning,
      plotmarks, angles, quotes
    }
    \tikzset{>=latex} % sets default arrow type in tikz
    \tikzset{font=\small}
    
      
    \usepackage{multirow, bigdelim}
    
    % change to \usepackage[sols]{handout} to see solutions
    \usepackage[]{handout}
        
    \newcommand{\lheq}{\stackrel{\text{L'H}}{=}}
\newcommand{\nolheq}{\mathrel{\makebox[\widthof{$\lheq$}]{=}}}

    \definecolor{skyblue}{rgb}{0, 0.5, 1}
    \usepackage{tcolorbox}
\newenvironment{feedback}{%
  \begin{tcolorbox}[colframe=skyblue, colback=skyblue!10!white]
  }{\end{tcolorbox}}
      
    
    \newcommand\iscurrentenumi[1]{%
      \ifthenelse{\equal{#1}{\theenumi}}{}{#1}}
    \setenumerate[1]{ref=\#\arabic*}
    \setenumerate[2]{ref=\protect\iscurrentenumi{\theenumi}(\alph*)}

    
      
    
    \newlength{\tipmargin}
    \makeatletter
    
    \makeatother
    
      
    \ifthenelse{\isundefined{\C}}{%
      \newcommand{\C}{\mathbb{C}}%
    }{\renewcommand{\C}{\mathbb{C}}}
    \ifthenelse{\isundefined{\R}}{%
      \newcommand{\R}{\mathbb{R}}%
    }{\renewcommand{\R}{\mathbb{R}}}
    
    \definecolor{purple}{rgb}{0.5,0,1.0}
    \pgfplotsset{holdot/.style={color=red,fill=white,only marks,mark=*}}
    \pgfplotsset{soldot/.style={color=black,fill=black,only marks,mark=*}}
\pgfplotsset{compat=1.13}
\usetikzlibrary{quotes,arrows.meta}

    \pgfplotsset{holdot/.style={color=red,fill=white,only marks,mark=*}}

\newcommand{\answerbox}[3]{
    \fbox{\begin{minipage}[t][#1][c]{#2}
        Final answer: #3
    \end{minipage}}
    }

    \newcommand{\answerboxd}[2]{
    \fbox{\begin{minipage}[t][#1][c]{#2}
        Final answer: $f'(x)=$
    \end{minipage}}
    }
    
\begin{document}

 

 \noindent
{\large \mbox{} \hfill \bf Name:\underline{\hspace{3in}}}

 {\mbox{} \hfill \it (Please write your name neatly in print.)}


 
\medskip
 
\noindent
{\large\bf
Practice Math Ma Skills Check 2\\ Math Mb Spring 2025}
\noindent






\iffalse
\begin{center}
\begin{tabular}{|c|c|c|}
\hline
Problem & Points & Score\\ \hline
& &\\
1 &  & \\ 
\hline
&& \\
2&  &  \\ 
\hline
& &\\
3&&  \\ 
 \hline
& &\\
Total&  &  \\ 
& &\\ \hline

\end{tabular}
\end{center}
\fi


%{\large \em Please circle your section:}


%\begin{center}
%{\small
%\begin{tabular}{ccccccccc}
%Erica Dinkins & Matt Cavallo &  Shanelle Davis  & Robin Gottlieb & 	Mike Koske & Hakim Walker \\
%MWF 9am  & MWF 9am &  MWF 10:30am  &MWF 10:30am & MWF 12pm   & MWF 12pm\\
%\\\\
% &&  &     \\
% &&  &
%\end{tabular}

%}
%\vspace{ -.5in} 
%{\small
%\begin{tabular}{cccccccc}
 %Hannah Constantin & Matt Cavallo & 	Justin Hancock  & Grant Barkley \\
 %MWF 12pm &  MWF 12pm  & MWF 1:30pm  & MWF 3pm  \\
%\\\\
 %&&       \\
 %&&  
%\end{tabular}


%}
%\end{center}

{\Large
\begin{center}\textbf{Directions---Please read carefully!} \end{center}}

This is a practice test. The same directions will appear on the actual Skills Check, so read them carefully.

\hrulefill 

You have 60 minutes to complete this Skills Check. \textbf{Calculators or any other aids are not permitted.}  


For each problem, write your final answer in the answer box provided. You may do work on any page and may ask for scratch paper, but \textbf{no work or writing outside of the answer box will be graded.}

Write neatly, as answers that cannot be read will not receive credit. 

   

\begin{center}\textbf{Good luck!}\end{center}


\newpage

\section*{Derivatives}



\noindent Take the derivative of each of the following functions for Problems 1 through 5. You may use the rules of your choice. You do not need to simplify your answers.

\begin{enumerate}
    \item $f(x) = \sqrt{\ln(x-2)}$

    \pspace{\vfill}

    \begin{solution}
    We can use the Chain Rule. 
    \begin{align*}
    f'(x) &= \diff*{\left[\ln(x-2)\right]^{1/2}}{x} \\
    &= \frac{1}{2}\left[\ln(x-2)\right]^{-1/2} \cdot \diff*{\ln(x-2)}{x} \\
    &= \frac{1}{2}\left[\ln(x-2)\right]^{-1/2} \cdot \frac{1}{x-2} \cdot \diff*{(x-2)}{x} \\
    &= \frac{1}{2}\left[\ln(x-2)\right]^{-1/2} \cdot \frac{1}{x-2} \cdot 1
    \end{align*}

    \end{solution}
        
    \answerbox{0.8in}{5.5in}{$\displaystyle f'(x)=\solonly{\frac{1}{2}\left[\ln(x-2)\right]^{-1/2} \cdot \frac{1}{x-2}}$}

    \item $f(x) = \ln(\sqrt{2x-8})$

        \pspace{\vfill}

        \begin{solution}
            Before using the Chain Rule, we can use properties of logs to rewrite the function.
            \begin{align*}
                f(x) &= \ln\!\left[(2x-8)^{1/2}\right] \\
                &= \frac{1}{2}\ln(2x-8)
            \end{align*}
            Now we use the Chain Rule.
            \begin{align*}
                f'(x) &= \diff*{\frac{1}{2}\ln(2x-8)}{x} \\
                &= \frac{1}{2}\cdot \frac{1}{2x-8} \cdot \diff*{(2x-8)}{x}\\
                &= \frac{1}{2}\cdot \frac{1}{2x-8}\cdot 2
            \end{align*}
        \end{solution}

        \answerbox{0.8in}{5.5in}{$\displaystyle f'(x)=\solonly{\frac{1}{2x-8}}$}

\newpage 
    
    \item $f(x) = (e^{4x}+6)^5$

        \pspace{\vfill}

        \begin{solution}
            Chain Rule.
            \begin{align*}
                f'(x)&=\diff*{(e^{4x}+6)^5}{x} \\
                &= 5(e^{4x}+6)^4 \cdot \diff*{(e^{4x}+6)}{x} \\
                &= 5(e^{4x}+6)^4 \cdot 4e^{4x} \\ 
            \end{align*}
        \end{solution}

    \answerbox{0.8in}{5.5in}{$\displaystyle f'(x)=\solonly{20e^{4x}(e^{4x}+6)^4}$}

    \item $f(x) = (2+3^x) \cdot \ln x$ 

        \pspace{\vfill}
        \begin{solution}
            We can use the Product Rule.
            \begin{align*}
                f'(x) &= \diff*{\left[(2+3^x)\cdot \ln x\right]}{x}\\
                &= \left[\diff*{(2+3^x)}{x}\right] \cdot \ln x + (2+3^x) \cdot \diff*{\ln x}{x} \\
                &= 3^x \ln(3) \cdot \ln x + \frac{2+3^x}{x}
            \end{align*}
        \end{solution}
        
    \answerbox{0.8in}{5.5in}{$\displaystyle f'(x)=\solonly{3^x \ln(3) \ln x + \frac{2+3^x}{x}}$}

    \item $f(x) = \dfrac{e^{2x}+3}{4x^3}$

        \pspace{\vfill}

    \answerbox{0.8in}{5.5in}{$f'(x)=$}

\newpage 

    \item Find the equation of the line tangent to the graph of $y=4x^5-3x+7$ when $x=-1$.

        \vfill

        \answerbox{0.8in}{5.5in}{}
        
\end{enumerate}


\section*{Logs and Exponent Rules}

\begin{enumerate}[resume]
    \item Which of the following statements are true? Circle all correct answers.

    \begin{multicols}{2}
    \begin{enumerate}
        \item $e^{4\ln x}=4x$
        \item $\ln(x+y)=\ln x + \ln y$
        \item $e^{2\ln x}=x^2$
        \item $\dfrac{4a^{-2}}{b}=\dfrac{1}{4a^2b}$
        \item $\ln(xy^2)=\ln x + 2\ln(y)$
        \item $\ln(xy^2)=2\ln(xy)$
    \end{enumerate}
    \end{multicols}

\newpage 
    
    \item Simplify the following:
    \begin{multicols}{2}
            \begin{enumerate}
        \item $e^{\ln[(a-2b)^3]}$
        %\vfill 
       %     \answerbox{0.8in}{2.5in}
        \item $\ln\Big(\dfrac{4}{e^{6x-1}}\Big)$
       % \vfill
        %    \answerbox{0.8in}{2.5in}
    \end{enumerate}
    \end{multicols}

    \vfill 

        \begin{multicols}{2}
        \answerbox{0.8in}{2.5in}{}
                
        \answerbox{0.8in}{2.5in}{}
        
    \end{multicols}



\item Solve the following equations for $x$. Your final answer should not include any logarithms.

    \begin{multicols}{2}
            \begin{enumerate}
        \item $2\ln(3) - \ln(x) = \ln(x+6)$
        \item  $3^{2(x-2)}=81$
    \end{enumerate}
    \end{multicols}
\vfill 

    \begin{multicols}{2}
        \answerbox{0.8in}{2.5in}{}
                
        \answerbox{0.8in}{2.5in}{}
        
    \end{multicols}
    
\end{enumerate}


\newpage 

\section*{Limits}

Evaluate the following limits, if they exist. If the limit does not exist, write $\infty$, $-\infty$, or DNE where appropriate.

\begin{enumerate}[resume]

        
    \item $\displaystyle \lim_{x \to \infty} \dfrac{3x^2-2x+5}{4x^3}$

        \vfill

        \answerbox{0.8in}{5.5in}{}

        \item $\displaystyle \lim_{x \to  1} \dfrac{x^2+2x-3}{x^2-x}$

        \vfill

        \answerbox{0.8in}{5.5in}{}

\newpage 
        
    \item $\displaystyle \lim_{x \to  \infty} \dfrac{e^x + 1}{x^{101}}$

        \vfill

        \answerbox{0.8in}{5.5in}{}
        


        \item $\displaystyle \lim_{x \to 4} \dfrac{x^2-7}{x+6}$

        \vfill
        \answerbox{0.8in}{5.5in}{}
        
    \item $\displaystyle \lim_{x \to -2^{-}} \dfrac{3x-7}{x+2}$

        \vfill

        \answerbox{0.8in}{5.5in}{}
\end{enumerate}



\end{document}